% SPDX-License-Identifier: CC-BY-4.0
% Copyright (c) 2026 Martin Schröder <info@swedishembedded.com>

\section{Introduction}\label{sec:intro}

\paragraph{Question.}
Can a voice agent learn how to speak from the corrections its user gives it in conversation, so that
a corrected mistake stays corrected, nothing else gets worse, and the user can take a correction
back? This paper formalises that setting, specifies the experiments that would answer it, and
reports the evidence that exists. A finding belongs here if it changes what is claimed about, or how
one would measure, an agent that turns spoken or typed feedback about its own speech into lessons,
applies them, trains on them and releases the result.

\paragraph{Status of the claims.}
\Cref{tab:claims} lists the claims. At the time of this draft only one has partial evidence (the
latency of the speaking agent), one has failed its gate (speech understanding without a transcript),
and the rest are pending. The tables of this paper are generated from stored run reports
(\cref{sec:repro}); a claim cannot be marked met without a run that supports it.

\begin{table}[t]
\centering\small
\input{generated/claims_table}
\caption{The claims of this paper and the status of their evidence. ``Pending'': no test has run.
``Partial'': some of the gate is measured. ``Failed'': the test ran and the gate was not met.}
\label{tab:claims}
\end{table}

\paragraph{What is not claimed.}
Every voice in this work is a synthetic theatrical portrayal; no recording of Samuel Adams or Thomas
Jefferson exists, so no voice here is theirs, and every report and spoken session carries that label.
Until a real-speech test set is recorded and measured, nothing is claimed about real human speech.
``Trained persona'' is not claimed: the persona is a language model under a system prompt (\cref{sec:system}).
